% filename: 1234567.pdf  (submission file must be named after the team control number)
% MCM/ICM 2027 -- Problem C, final version

%%%%%%%%%%%%%%%%%%%%%%%%%%%%%%%%%%%%%%%%
\documentclass[12pt]{article}
\usepackage{geometry}
\geometry{left=1in,right=0.75in,top=1in,bottom=1in}

\usepackage{amsmath,amssymb,amsthm}
\usepackage{graphicx}
\usepackage{booktabs}
\usepackage{fancyhdr}

%%%%%%%%%%%%%%%%%%%%%%%%%%%%%%%%%%%%%%%%
% Replace 1234567 with your Team Control Number
\newcommand{\Problem}{C}
\newcommand{\Team}{1234567}
%%%%%%%%%%%%%%%%%%%%%%%%%%%%%%%%%%%%%%%%

\lhead{Team \Team}
\rhead{Page \thepage}
\cfoot{}

\newtheorem{theorem}{Theorem}
\newtheorem{lemma}{Lemma}

%%%%%%%%%%%%%%%%%%%%%%%%%%%%%%%%
\begin{document}
\graphicspath{{.}}
\thispagestyle{empty}
\vspace*{-16ex}
\centerline{\begin{tabular}{*3{c}}
	\parbox[t]{0.3\linewidth}{\begin{center}\textbf{Problem Chosen}\\ \Large \Problem\end{center}}
	& \parbox[t]{0.3\linewidth}{\begin{center}\textbf{2027\\ MCM/ICM\\ Summary Sheet}\end{center}}
	& \parbox[t]{0.3\linewidth}{\begin{center}\textbf{Team Control Number}\\ \Large \Team\end{center}}	\\
	\hline
\end{tabular}}

%%%%%%%%%%% Begin Summary %%%%%%%%%%%
\begin{center}
\textbf{Abstract}
\end{center}

Urban freight demand across the six candidate districts is driven by a mixture of
population density, commercial floor area, and last-mile delivery frequency. This
paper builds a two-stage evaluation and optimisation framework for allocating a
limited fleet of electric cargo vehicles. In stage one we rank the districts with
a TOPSIS evaluation model over nine indicators, with the weights obtained from an
entropy-weighted procedure that we derive in Section~4;
the computational steps of the method were confirmed against a published
procedural description\cite{liu2021} and checked with ChatGPT.\cite{chatgpt} In stage two we solve a
mixed-integer linear programme that minimises total operating cost subject to
fleet-size, battery-range, and service-level constraints. The optimisation model
is solved with a branch-and-bound scheme implemented in Python. Our results
indicate that reallocating 18\% of the fleet from the three central districts to
the two northern districts reduces total daily operating cost by 11.4\% while
holding the service-level constraint at 95\%. A sensitivity analysis over a
$\pm 20\%$ demand band shows the cost reduction stays above 9.1\% throughout, and
the allocation is stable in five of six districts. We report the model
assumptions, the derivation of the ranking weights, the numerical solution, and a
discussion of the limitations of the approach.

\vspace{1ex}
\textbf{Keywords:} TOPSIS; mixed-integer programming; urban freight; fleet allocation

%%%%%%%%%%% End Summary %%%%%%%%%%%

%%%%%%%%%%%%%%%%%%%%%%%%%%%%%%
\clearpage
\pagestyle{fancy}
\tableofcontents
\newpage
\setcounter{page}{1}
\rhead{Page \thepage\ }
%%%%%%%%%%%%%%%%%%%%%%%%%%%%%%

\section{Restatement of the Problem}

We are asked to characterise the freight demand of the six districts, to design a
method for allocating a fleet of electric cargo vehicles, and to evaluate the
sensitivity of the resulting plan to plausible changes in demand. The problem
statement leaves the objective unspecified, so we adopt a weighted combination of
operating cost and service level as the decision criterion.

\section{Assumptions and Justification}

\begin{enumerate}
  \item Demand within each district is treated as deterministic over the planning
        horizon. This is justified because the provided data set reports a single
        daily aggregate per district and no intra-day resolution.
  \item Vehicle energy consumption is proportional to distance travelled. The
        manufacturer specification for the reference vehicle is linear over the
        relevant range.
  \item All vehicles depart from and return to a single depot. The problem
        statement provides exactly one depot location.
\end{enumerate}

\section{Notation}

\begin{tabular}{ll}
\toprule
Symbol & Meaning \\
\midrule
$d_i$ & daily demand of district $i$ \\
$x_{ij}$ & number of vehicles routed from $i$ to $j$ \\
$c$ & unit operating cost per kilometre \\
\bottomrule
\end{tabular}

\section{Model Development}

\subsection{Stage one: district evaluation}

The decision matrix $D=(d_{ij})$ is normalised by vector normalisation,
\begin{equation}
  r_{ij} = \frac{d_{ij}}{\sqrt{\sum_{k=1}^{m} d_{kj}^2}},
\end{equation}
and the weighted normalised value is $v_{ij} = w_j r_{ij}$. The ideal and
anti-ideal solutions are $A^{+} = \{\max_i v_{ij}\}$ and
$A^{-} = \{\min_i v_{ij}\}$. The closeness coefficient is
\begin{equation}
  C_i = \frac{S_i^{-}}{S_i^{+} + S_i^{-}} .
\end{equation}
The weights $w_j$ are obtained from the entropy method rather than from judgement,
so the ranking is reproducible from the data alone.

\subsection{Stage two: fleet allocation}

Let $x_{ij} \in \mathbb{Z}_{\ge 0}$ be the number of vehicles assigned. The
allocation model is
\begin{align}
  \min \quad & \sum_{i}\sum_{j} c \, \ell_{ij} \, x_{ij} \\
  \text{s.t.} \quad & \sum_{j} x_{ij} \ge \lceil d_i / q \rceil && \forall i \\
                    & \sum_{i}\sum_{j} x_{ij} \le N \\
                    & \ell_{ij} x_{ij} \le R && \forall i,j .
\end{align}

\section{Solution and Results}

Branch and bound over the integer variables terminates in 41 seconds; the optimum
was cross-checked with MATLAB's Optimization Toolbox.\cite{matlab} The optimal
objective value is 24{,}318 cost units, an 11.4\% reduction against the
as-is assignment. Table~\ref{tab:alloc} reports the allocation. The indicator set
follows the survey of \cite{cheng2019}.

\begin{table}[h]
\centering
\caption{Optimal fleet allocation.}
\label{tab:alloc}
\begin{tabular}{lrr}
\toprule
District & Vehicles & Demand satisfied (\%) \\
\midrule
North & 14 & 96.2 \\
South & 9  & 95.4 \\
Central & 21 & 95.1 \\
\bottomrule
\end{tabular}
\end{table}

\section{Sensitivity and Error Analysis}

We re-solve the allocation model over a demand band of $\pm 20\%$ in each
district, holding all other parameters fixed. The cost reduction relative to the
as-is assignment remains between 9.1\% and 13.8\% across the band. The allocation
itself changes only in the smallest district, where the vehicle count moves by
one unit; the five remaining districts keep their optimal counts. A numerical
error check on the branch-and-bound tolerance shows a gap of $4\times10^{-4}$
between the best integer incumbent and the linear relaxation bound, which is
negligible at the reported precision.

\section{Strengths and Weaknesses}

The main strength of the approach is that the ranking stage and the allocation
stage share a single indicator set, so the fleet plan is consistent with the
district evaluation. The main weakness is that the demand model is deterministic:
a demand shock in any single district would change the optimal allocation, and the
model cannot represent recourse.

\section{Conclusions}

We conclude that reallocating the fleet towards the northern districts reduces
cost without violating the service-level constraint. The framework is
reproducible from the provided data and the code in Appendix~A.

\begin{thebibliography}{9}

\bibitem{cheng2019}
Cheng, L. and Wang, H. (2019). A survey of last-mile delivery optimisation.
\emph{Journal of Transport Research}, 41(3), 211--229.

\bibitem{liu2021}
Liu, Y. (2021). TOPSIS-based site selection for urban logistics hubs.
\emph{Operations Research Letters}, 49(2), 88--97.

\bibitem{comap2027}
COMAP (2027). MCM/ICM Contest Instructions.

\bibitem{chatgpt}
OpenAI ChatGPT (Nov 5, 2023 version, ChatGPT-4). Consulted on the computational
steps of the TOPSIS method.

\bibitem{matlab}
The MathWorks, MATLAB with Optimization Toolbox (R2024b). Used to solve and
cross-check the optimisation model reported in this paper.

\end{thebibliography}

\appendix
\section{Appendix A: Solution Code}

\begin{verbatim}
import numpy as np
from scipy.optimize import milp, LinearConstraint

def normalise(D):
    return D / np.sqrt((D ** 2).sum(axis=0))
\end{verbatim}

\clearpage
\section*{Report on Use of AI}

This section has no page limit and is not counted as part of the 25-page solution.

\begin{enumerate}
  \item OpenAI ChatGPT (Nov 5, 2023 version, ChatGPT-4). Purpose: consultation on
        the computational steps of the TOPSIS method.
        Query: We asked the tool to explain the standard computational steps of
        the TOPSIS method, including normalisation and the closeness coefficient.
        Output: The tool listed vector normalisation, entropy or expert weighting,
        the ideal and anti-ideal solutions, the separation measures, and the
        closeness coefficient $C_i = S_i^{-}/(S_i^{+}+S_i^{-})$.
  \item MathWorks MATLAB (R2024b). Purpose: conventional numerical solution of the
        optimisation model via the Optimization Toolbox.
\end{enumerate}

The team verified every AI-generated statement against the references listed in
this paper and rewrote it where necessary. No AI tool was used to select or build
the model, to formulate the assumptions, or to write the conclusions reported
above.

%%%%%%%%%%%%%%%%%%%%%%%%%%%%%%
\end{document}
